\section{Related Work}

% Progress Report: one short paragraph for 2.1, 2.3, 2.4 and 2.5, one sentence for 2.2 (1 to 1.5 pages).
% Each subsection must justify one choice of the methodology. Write it as an argument, not as a list.
% Candidate works are in literature.bib; every entry there is UNVERIFIED and must be checked
% against the original paper before it is cited.

\subsection{Hand Detection and Landmark Estimation}
% JUSTIFIES: why MediaPipe Hands.
% COVER: classical methods (skin colour, contour, hand-crafted features) and their sensitivity to
% lighting and background; MediaPipe Hands, a palm detector followed by a 21-landmark regression model.
% CITE (after checking): zhang2020mediapipe
\todo{2.1}

\subsection{Gesture Datasets and Benchmarks}
% JUSTIFIES: why we record our own data.
% COVER: size, classes, static or dynamic content and fit for device control of public datasets;
% none of them is recorded under controlled, labelled lighting, background and distance conditions.
% CITE (after checking): kapitanov2022hagrid, materzynska2019jester, benitezgarcia2020ipn
\todo{2.2 (Progress Report: one sentence)}

\subsection{Static and Dynamic Gesture Recognition}
% JUSTIFIES: why landmark-based features and small models.
% COVER: landmark-based versus video-based recognition, classical versus learned classifiers.
% CITE (after checking): molchanov2016online
\todo{2.3}

\subsection{Continuous Recognition and False-Trigger Prevention}
% JUSTIFIES: why segmentation with hysteresis, a hold time and a cooldown.
% COVER: how existing work finds the start and end of a gesture, rejects movements that are not
% gestures, and measures false positives and latency.
% CITE (after checking): alon2009unified, jacob1990what
\todo{2.4}

\subsection{Temporal Smoothing}
% JUSTIFIES: why the 1-euro filter.
% COVER: noisy landmark streams, the jitter versus lag trade-off.
% CITE (after checking): casiez2012oneeuro
\todo{2.5}

\iffinal
\subsection{Comparison and Positioning}
% WRITE: a table of the cited works and one sentence on how this project differs, for example:
% existing work mostly reports accuracy under one condition, while this project evaluates
% robustness across users, backgrounds, lighting and distance, and measures false-trigger
% rate and latency explicitly.
\todo{comparison table and positioning}
\fi
